% Generated draft: requires mathematical and source review.
\documentclass[11pt]{article}
\usepackage[margin=25mm]{geometry}
\usepackage{fontspec}
\setmainfont{Latin Modern Roman}
\usepackage{amsmath,amssymb,mathtools,unicode-math}
\setmathfont{Latin Modern Math}
\usepackage{enumitem,needspace,xurl,hyperref}
\setlength{\parindent}{0pt}
\setlength{\parskip}{5pt}
\setlength{\emergencystretch}{3em}
\newcommand{\N}{\mathbb{N}}
\newcommand{\Z}{\mathbb{Z}}
\newcommand{\Q}{\mathbb{Q}}
\newcommand{\R}{\mathbb{R}}
\newcommand{\C}{\mathbb{C}}
\newcommand{\F}{\mathbb{F}}
\newcommand{\A}{\mathbb{A}}
\newcommand{\PP}{\mathbb{P}}
\newcommand{\E}{\mathbb{E}}
\newcommand{\eps}{\varepsilon}
\newcommand{\dd}{\,\mathrm{d}}
\begin{document}
\section*{UnsolvedMath ID: 30001601}
\textbf{Source problem code:} catalogue entry
\subsection{Definitions and mathematical statement}
Let \(S_1,\ldots,S_r\) be a finite collection of equisingular types of plane-curve singularities, as in the supplied record. Write \(p_{ij}\) for the associated proper or infinitely near points appearing in its blowup construction, and put \(\mu_{ij}=\operatorname{mult}_{p_{ij}}S_i\). The record does not specify the precise notion of equisingular type or which points of a resolution or cluster are included. Let \(m_1\geq m_2\geq m_3\) denote the three largest associated multiplicities in the sense intended by the source; the collection from which they are selected is not specified.

Let \(\pi:X\to\mathbb{P}^{2}\) be a general proximity-compatible blowup realizing the prescribed data. Here \(L\) denotes the class of a line in \(\mathbb{P}^{2}\), and \(E_{ij}\) denotes the exceptional divisor class assigned to \(p_{ij}\) in the source formula. The record leaves unspecified the base field, the parameter space underlying ``general,'' the required proximity relations, and whether \(E_{ij}\) means a strict transform or a total transform. For an integral divisor class \(D\) on a specified \(X\), let \(|D|\) denote its complete linear system, with \(\dim|D|=h^0(X,\mathcal{O}_X(D))-1\) when nonempty and \(\dim|D|=-1\) when empty.

\textbf{Conjectural equality supplied by the record.} For every collection of types and general proximity-compatible blowup as above, and every integer \(d\geq m_1+m_2+m_3\), is it true that
\[
\dim\left|d\pi^{*}L-\sum_{i,j}\mu_{ij}E_{ij}\right|
=
\max\left\{-1,\binom{d+2}{2}-C(S_1,\ldots,S_r)-1\right\}?
\]
The quantity \(C(S_1,\ldots,S_r)\) is not defined in the supplied record. No formula for it is assumed here. Consequently, this is a faithful transcription of the intended equality with explicit unresolved qualifications, rather than a fully specified mathematical conjecture. A precise formulation requires a definition of \(C\), the cluster and multiplicity conventions, the base field, and the meaning of generality. The record also gives no convention for cases with fewer than three associated multiplicities.
\subsection{Short English statement}
The proposed equality says that, once the degree is at least the sum of the three largest associated multiplicities, the prescribed plane-curve system has the dimension displayed in the source formula. The invariant \(C\) and several geometric conventions remain undefined. Supplied dated curation labels the problem partially solved, reporting the ADE case proved and arbitrary singularity types still open; these claims are not independently verified here.
\subsection{Sources}
\begin{itemize}
\item Supplied local UnsolvedMath registry record: permanent ID \(30001601\), source problem code \texttt{catalogue entry}, title ``Expected Dimensions of Equisingular Plane-Curve Systems,'' in the Oberwolfach Reports Open Problems collection. The record gives proposed year 2010 and no named proposer.
\item Supplied source citation: \emph{Mini-Workshop: Linear Series on Algebraic Varieties} (2011), \emph{Oberwolfach Reports}, pp.\ 2613--2650, DOI \texttt{10.4171/owr/2010/45}. Link: \url{https://doi.org/10.4171/owr/2010/45}. The supplied background describes Oberwolfach Reports as documenting mathematical workshops and their open-problem discussions.
\item Supplied literature source: \url{https://arxiv.org/abs/math/0612310}. No title, authors, or claim-specific passage is supplied for this link.
\item Supplied literature source: \url{https://arxiv.org/abs/2008.02640}. No title, authors, or claim-specific passage is supplied for this link.
\item Literature assessment dated exactly \texttt{2026-08-21}: the supplied curation attributes the conjectural dimension formula to Greuel--Lossen--Shustin and states that it is proved for ADE singularities under the stated degree hypothesis but remains open for arbitrary singularity types. This is supplied dated curation, not an independently verified literature assessment. Its background expressly requests independent verification of time-sensitive status and cited claims.
\end{itemize}
\end{document}
